\documentclass[
  spanish,
  aspectratio=169,
  xcolor={dvipsnames,table}
]{beamer}

\geometry{paperwidth=19.2cm,paperheight=10.8cm}
\usepackage[utf8]{inputenc}
\usepackage[T1]{fontenc}
\usepackage[spanish,es-tabla]{babel}
\usepackage[scaled=.96]{helvet}
\renewcommand{\familydefault}{\sfdefault}
\usepackage{booktabs}
\usepackage{graphicx}
\usepackage{ragged2e}
\usepackage{tikz}
\usepackage{hyperref}
\usetikzlibrary{calc,positioning,fit,shadows.blur}

\newcommand{\studentname}{Martin Jonathan de la Cruz}
\newcommand{\studentshort}{M. J. de la Cruz}
\newcommand{\studentid}{ES2611202040}
\newcommand{\universityname}{Universidad Abierta y a Distancia de México}
\newcommand{\facultyname}{Licenciatura en Derecho}
\newcommand{\coursename}{Garantías Constitucionales}
\newcommand{\coursecode}{LDE-S2B1}
\newcommand{\activitytitle}{Bienvenida e introducción a Garantías Constitucionales}
\newcommand{\activitysubtitle}{Objetivo, diagnóstico y ruta inicial de aprendizaje}
\newcommand{\activityweek}{Inicio de asignatura}
\newcommand{\teachingfigure}{Aarón López Pérez}
\newcommand{\deliverydate}{\today}
\newcommand{\locationname}{Roma Norte, Ciudad de México}
\newcommand{\departmentlogo}{img/departamentos/UnADM.pdf}

\definecolor{unadmGreenDark}{HTML}{174A3A}
\definecolor{unadmGreen}{HTML}{5F8F3A}
\definecolor{unadmGold}{HTML}{B88A2A}
\definecolor{unadmPaper}{HTML}{F6F7F2}
\definecolor{unadmInk}{HTML}{1F2A24}

\mode<presentation>{
  \usetheme{default}
  \usefonttheme{professionalfonts}
  \setbeamertemplate{navigation symbols}{}
  \setbeamertemplate{blocks}[rounded][shadow=false]
}
\setbeamersize{text margin left=0.60cm,text margin right=0.60cm}
\setbeamercolor{background canvas}{bg=white}
\setbeamercolor{normal text}{fg=unadmInk,bg=white}
\setbeamercolor{structure}{fg=unadmGreenDark}
\setbeamercolor{frametitle}{fg=white,bg=unadmGreenDark}
\setbeamercolor{block title}{fg=white,bg=unadmGreen}
\setbeamercolor{block body}{fg=unadmInk,bg=unadmPaper}
\setbeamerfont{title}{size=\LARGE,series=\bfseries}
\setbeamerfont{frametitle}{size=\large,series=\bfseries}
\setbeamertemplate{itemize item}{\textcolor{unadmGreen}{\large$\blacktriangleright$}}

\setbeamertemplate{frametitle}{
  \nointerlineskip
  \begin{beamercolorbox}[wd=\textwidth,ht=1.02cm,dp=0.22cm,leftskip=0.35cm,rightskip=0.25cm]{frametitle}
    \insertframetitle\hfill\includegraphics[height=0.60cm]{\departmentlogo}
  \end{beamercolorbox}
  {\color{unadmGold}\rule{\textwidth}{1.3pt}}
}

\setbeamertemplate{footline}{
  \leavevmode
  \hbox{
    \begin{beamercolorbox}[wd=.34\paperwidth,ht=2.7ex,dp=1ex,leftskip=1em]{author in head/foot}\color{white}\studentshort\end{beamercolorbox}
    \begin{beamercolorbox}[wd=.40\paperwidth,ht=2.7ex,dp=1ex,center]{title in head/foot}\color{white}\coursename\end{beamercolorbox}
    \begin{beamercolorbox}[wd=.26\paperwidth,ht=2.7ex,dp=1ex,rightskip=1em plus 1fil]{date in head/foot}\color{white}\coursecode\hfill\insertframenumber/\inserttotalframenumber\end{beamercolorbox}
  }
}
\setbeamercolor{author in head/foot}{bg=unadmGreenDark}
\setbeamercolor{title in head/foot}{bg=unadmGreen}
\setbeamercolor{date in head/foot}{bg=unadmGreenDark}

\title[\coursecode]{\activitytitle}
\subtitle{\activitysubtitle}
\author[\studentshort]{\studentname\\\footnotesize Matricula: \studentid}
\institute[UnADM]{\universityname\\\facultyname}
\date[\activityweek]{\locationname\\\activityweek\\\deliverydate}

\begin{document}

\begin{frame}[plain]
  \begin{tikzpicture}[remember picture,overlay]
    \fill[unadmGreenDark] (current page.south west) rectangle ([xshift=0.58\paperwidth]current page.north west);
    \fill[unadmGreen] ([xshift=0.58\paperwidth]current page.south west) rectangle (current page.north east);
    \node[anchor=center,opacity=0.10] at (current page.center) {\includegraphics[width=0.72\paperwidth]{\departmentlogo}};
    \node[anchor=north west] at ([xshift=0.58cm,yshift=-0.46cm]current page.north west) {\includegraphics[width=2.55cm]{\departmentlogo}};
  \end{tikzpicture}
  \vspace*{1.62cm}
  {\color{white}\fontsize{24}{29}\selectfont\bfseries \activitytitle\par}
  \vspace{0.28cm}
  {\color{white!86}\large \activitysubtitle}\\[0.35cm]
  {\color{unadmGold}\rule{0.58\linewidth}{1.3pt}}\\[0.35cm]
  {\color{white}\studentname}\\
  {\color{white!82}\footnotesize \facultyname}
\end{frame}

\begin{frame}{Mensaje de bienvenida}
  \begin{block}{Inicio de la asignatura}
    Iniciaremos la asignatura de \textbf{Garantías Constitucionales}. Sean todas y todos bienvenidos.
  \end{block}
  \begin{itemize}
    \item La sesión sincrónica permitirá interactuar y conocerse.
    \item Los comentarios e inquietudes serán escuchados.
    \item El punto de partida será un marco de respeto.
  \end{itemize}
\end{frame}

\begin{frame}{Presentación del docente}
  \begin{block}{Aarón López Pérez}
    Abogado con posgrados de Especialidad y Maestría realizados en la Facultad de Derecho de la Universidad Nacional Autónoma de México.
  \end{block}
  \begin{itemize}
    \item Propone aprender conjuntamente a partir de las opiniones del grupo.
    \item Establece un tono cercano, académico e institucional.
  \end{itemize}
\end{frame}

\begin{frame}{Introducción y objetivo}
  \begin{block}{Valor formativo}
    La asignatura es fundamental para comprender y ejercer la abogacía porque conecta derechos de las personas, defensa y protección constitucional.
  \end{block}
  \begin{itemize}
    \item Fundamento: Constitución mexicana.
    \item Objeto: derechos humanos, garantías y medios de control constitucional.
    \item Finalidad: defensa profesional de derechos, justicia y desarrollo social.
  \end{itemize}
\end{frame}

\begin{frame}{Conceptos clave}
  \begin{tabular}{p{0.34\linewidth}p{0.56\linewidth}}
    \toprule
    \textbf{Concepto} & \textbf{Sentido inicial} \\
    \midrule
    Garantías Constitucionales & Derechos de las personas y mecanismos para su defensa y protección. \\
    Derechos humanos & Núcleo de análisis y defensa de la asignatura. \\
    Medios de control constitucional & Vías previstas en la Constitución para proteger derechos. \\
    Diagnóstico inicial & Base para retroalimentar y nivelar pensamiento jurídico. \\
    \bottomrule
  \end{tabular}
\end{frame}

\begin{frame}{Ruta de aprendizaje}
  \begin{enumerate}
    \item Diagnosticar conocimientos previos con microcasos.
    \item Distinguir derecho, garantía y medio de control.
    \item Identificar hechos, autoridad o acto, derecho afectado y remedio posible.
    \item Validar la respuesta con fundamento constitucional.
  \end{enumerate}
\end{frame}

\begin{frame}{Términos relacionados}
  \begin{columns}[T]
    \begin{column}{0.48\linewidth}
      \begin{block}{Eje humano}
        \begin{itemize}
          \item Bienvenida.
          \item Estudiantes.
          \item Presentación del docente.
        \end{itemize}
      \end{block}
    \end{column}
    \begin{column}{0.48\linewidth}
      \begin{block}{Eje jurídico}
        \begin{itemize}
          \item Abogacía.
          \item Constitución mexicana.
          \item Defensa de derechos humanos.
          \item Justicia y desarrollo social.
        \end{itemize}
      \end{block}
    \end{column}
  \end{columns}
\end{frame}

\begin{frame}{Cierre}
  \begin{block}{Pauta de trabajo}
    La asignatura inicia con una invitación a participar, escuchar y construir aprendizaje jurídico orientado a la protección constitucional de derechos.
  \end{block}
  \begin{itemize}
    \item Tono: respetuoso.
    \item Dinámica: sesión sincrónica y diagnóstico inicial.
    \item Propósito: aprendizaje conjunto y defensa profesional de derechos humanos.
  \end{itemize}
\end{frame}

\end{document}